% ---- pg-505 （书页 493） ----

体结论, 要求有一个公共的协方差矩阵 $\boldsymbol{\Sigma}$. 若这一点是基本真实的, 且样本足够大, 则法则(11-67)应该表现相当好. 在任何情况下, 其性能可通过计算估计的失误率来检验. 具体来说, 应该计算由\hyperref[eq:11-57]{式(11-57)}给出的期望真实失误率的拉亨布鲁克估计值.

\section{逻辑斯蒂回归与分类}
\phantomsection\label{sec:11.8}

\subsection{引言}

前面讨论的分类函数是基于定量变量. 这里我们讨论一种分类方法, 其中某些或所有变量是定性的. 此方法称为逻辑斯蒂(logistic)回归. 在最简单的情况, 响应变量 $Y$ 只取两个值. 例如, $Y$ 可能是“男性”或“女性”, 或者“就业”或“失业”.

即使响应可能是两值的定性变量, 我们总可以将两个值记为 0 和 1. 例如, 可以令“男性”$=0$, “女性”$=1$. 则 1 的概率 $p$ 是我们感兴趣的参数. 此概率表示总体中标识为 1 的人所占比例. 分布的均值也是 $p$, 因为
\[
\text{均值}=0\times(1-p)+1\times p=p
\]
0 的比例为 $1-p$, 常记为 $q$.

分布的方差为
\[
\text{方差}=0^2\times(1-p)+1^2\times p-p^2=p(1-p)
\]
显然方差不是常数. 若 $p=0.5$, 则方差$=0.5\times 0.5=0.25$; 而当 $p=0.8$, 则方差$=0.8\times 0.2=0.16$. 当 $p$ 趋于 0 或 1 时, 方差趋于 0.

令响应 $Y$ 取值 0 或 1. 若我们要对取 1 的概率建立有一个预测变量的线性模型, 我们可以令
\[
p=E(Y\mid z)=\beta_0+\beta_1 z
\]
并再加上一个误差项 $\varepsilon$. 但是这个模型有严重缺陷.

\begin{itemize}
\item 由于线性函数是无界的, 因此响应 $Y$ 的预测值可能大于 1 或小于 0.
\item 回归分析的一个假定是: $Y$ 的方差关于预测变量 $Z$ 的所有值保持不变. 上面已经说明此假定在此不适合. 当然, 如果用加权的最小二乘法可以对此有适当改进.
\end{itemize}

我们需要其他方法将预测变量或协变量 $Z$ 引入模型(见文献[26]). 自此以后, 我们总在协变量的观测 $Z=z$ 给定的条件下对 $p(z)$ 建模, 即使实际上协变量是不固定的.

\subsection{逻辑模型}

不是直接对概率 $p$ 建立模型, 我们先考虑“机会比”(odds ratio)
\[
\mathrm{odds}=\frac{p}{1-p}
\]
即 1 的概率对 0 的概率之比. 注意: 与概率不同, 机会比可以大于 1. 若 80\% 的旅客的行李不被海关检查, 则 $p=0.8$, 则不被检查的机会是 $0.8/0.2=4$, 或 4 对 1. 但是机会比是不对称的, 因为被检查的机会是 $0.2/0.8=1/4$. 取自然对数, 我们得到 $\ln(4)=1.386$, 而 $\ln(1/4)=-1.386$, 刚好相差一个符号.

考虑机会比的自然对数, 其图像在\hyperref[fig:11.15]{图 11.15}. 当机会 $x$ 为 1, 即出现 0 或 1 的机会相同, 则 $x$ 的自然对数为 0. 当机会 $x$ 大于 1, 则自然对数随 $x$ 的增加而缓慢增加. 然而, 当机会 $x$ 小于 1, 则自然对数随 $x$ 的下降而迅速增加到 0.

在对于二值变量的逻辑斯蒂回归模型中, 我们对机会比的自然对数(称为 $p$ 的“logit”)建模.
